\documentclass[a4paper, serif, 10pt]{moderncv}

%% moderncv
% color and style
\moderncvcolor{black}
\moderncvstyle{banking}

%% page layout/margins
\usepackage[top=2.5cm, bottom=2.5cm, left=1.5cm, right=1.5cm]{geometry}

\nopagenumbers{}

%% Babel config for Indonesian (Bahasa Indonesia) language
% ref:
% - https://latex3.github.io/babel/guides/locale-indonesian.html
\usepackage[indonesian]{babel}

%% fontspec
\usepackage{fontspec}

\IfFontExistsTF{Linux Libertine}{
  \setmainfont[Ligatures={TeX, Common}, Numbers=OldStyle]{Linux Libertine}
}{}
\IfFontExistsTF{Linux Libertine O}{
  \setmainfont[Ligatures={TeX, Common}, Numbers=OldStyle]{Linux Libertine O}
}{}

%% CTeX
% CJK support, used to show CJK characters in the resume
%
% - heading=false: disable ctex heading
% - punct=kaiming: use kaiming punctuations styles for CJK
% - scheme=plain: use plain scheme, do not override \normalsize font size
% - space=auto: space settings for CJK characters
%
% CJK fonts are resolved per locale language, e.g. Noto Serif/Sans CJK JP
% for Japanese resumes, with fallback to the next region if not installed
%
% ref:
% - http://ctan.mirrorcatalogs.com/language/chinese/ctex/ctex.pdf
\usepackage[UTF8, heading=false, punct=kaiming, scheme=plain, space=auto]{ctex}

\IfFontExistsTF{Noto Serif CJK SC}{
  \setCJKmainfont{Noto Serif CJK SC}
}{}
\IfFontExistsTF{Noto Sans CJK SC}{
  \setCJKsansfont{Noto Sans CJK SC}
}{}

%% line spacing
\usepackage{setspace}
\setstretch{1.125}

%% URL styling - use normal text instead of monospace
\urlstyle{same}

%% Patch \httplink and \httpslink to handle full URLs
% The original moderncv macros always prepend http:// or https:// to URLs.
% This causes issues when the URL already has a protocol (e.g., https://example.com)
% resulting in malformed URLs like https://https://example.com.
% This patch checks if the URL already contains "://" and uses it directly if so.
% Note: We use \str_set:Nx (x-expansion) to expand macros like \@homepage first.
\ExplSyntaxOn
\str_new:N \g_yamlresume_url_str

\RenewDocumentCommand{\httpslink}{O{}m}{
  \str_set:Nx \g_yamlresume_url_str {#2}
  \str_if_in:NnTF \g_yamlresume_url_str {://}
    { \url{#2} }
    { \url{https://#2} }
}

\RenewDocumentCommand{\httplink}{O{}m}{
  \str_set:Nx \g_yamlresume_url_str {#2}
  \str_if_in:NnTF \g_yamlresume_url_str {://}
    { \url{#2} }
    { \url{http://#2} }
}
\ExplSyntaxOff

\name{Budi Santoso}{}
\title{Manajer SDM Senior | Spesialis Hubungan Karyawan \& Rekrutmen}
\phone[mobile]{+62 812-3456-7890}
\email{budi.santoso@email.com}
\homepage{https://www.budisantoso-hr.id}

\address{Jl. Jenderal Sudirman Kav. 52-53, Jakarta Selatan, DKI Jakarta, Indonesia, 12190}{}{}

\extrainfo{{\small \faLinkedin}\ \href{https://www.linkedin.com/in/budisantoso-hr}{@budisantoso-hr} {} {} {} • {} {} {} 
{\small \faTwitter}\ \href{https://twitter.com/budisantoso\_hr}{@budisantoso\_hr}}

\begin{document}

\maketitle

\section{Informasi Dasar}

\cvline{}{\begin{itemize}
\item Manajer Sumber Daya Manusia dengan pengalaman lebih dari 10 tahun di perusahaan multinasional dan startup teknologi.
\item Ahli dalam merancang strategi rekrutmen, mengelola hubungan industrial, dan meningkatkan keterlibatan karyawan.
\item Berpengalaman memimpin tim HR lintas fungsi dan mengimplementasikan sistem HRIS untuk efisiensi operasional.
\item Komitmen tinggi terhadap pengembangan budaya kerja inklusif, berintegritas, dan berorientasi pada hasil.
\end{itemize}}
\section{Pendidikan}

\cventry{Sep 2008–Jul 2012}
        {Sarjana, Psikologi, Nilai: 3.75}
        {Universitas Indonesia}
        {\url{https://www.ui.ac.id}}
        {}
        {\begin{itemize}
\item Lulus dengan predikat cum laude dan meraih penghargaan mahasiswa berprestasi.
\item Aktif dalam organisasi kemahasiswaan sebagai ketua divisi pengembangan sumber daya anggota.
\end{itemize}
\textbf{Mata Kuliah}: Psikologi Industri dan Organisasi, Manajemen Sumber Daya Manusia, Perilaku Organisasi, Metodologi Penelitian, Psikometri, Konseling Karir}

\cventry{Agu 2013–Jun 2015}
        {Magister, Manajemen Sumber Daya Manusia, Nilai: 3.90}
        {Universitas Gadjah Mada}
        {\url{https://www.ugm.ac.id}}
        {}
        {\begin{itemize}
\item Tesis tentang pengaruh budaya organisasi terhadap keterlibatan karyawan di sektor teknologi.
\item Menerima beasiswa penuh dari program pascasarjana.
\end{itemize}}
\section{Pengalaman Kerja}

\cventry{Mar 2019–Sekarang}
        {Manajer Sumber Daya Manusia}
        {PT Maju Bersama Sejahtera}
        {\url{https://www.majubersama.co.id}}
        {}
        {\begin{itemize}
\item Memimpin tim HR dengan 8 anggota untuk mendukung 500+ karyawan di seluruh Indonesia.
\item Mengembangkan dan menerapkan strategi rekrutmen yang menurunkan time-to-hire sebesar 30\%.
\item Merancang program retensi karyawan yang meningkatkan engagement score dari 72\% menjadi 85\%.
\item Mengelola hubungan industrial, termasuk negosiasi perjanjian kerja bersama dan penanganan perselisihan.
\item Mengimplementasikan sistem HRIS berbasis cloud untuk otomatisasi payroll, absensi, dan manajemen kinerja.
\end{itemize}
\textbf{Kata Kunci}: Manajemen Talenta, Hubungan Industrial, HRIS, Engagement, Kepatuhan}

\cventry{Jun 2015–Feb 2019}
        {Supervisor Sumber Daya Manusia}
        {PT Karya Digital Nusantara}
        {\url{https://www.karyadigital.id}}
        {}
        {\begin{itemize}
\item Mengawasi proses rekrutmen end-to-end untuk posisi teknologi dan non-teknologi.
\item Mengembangkan program onboarding yang meningkatkan retensi karyawan baru sebesar 20\%.
\item Mengelola administrasi kompensasi dan benefit, termasuk BPJS Kesehatan dan Ketenagakerjaan.
\item Menjadi penghubung antara manajemen dan karyawan dalam menangani keluhan dan mediasi.
\end{itemize}
\textbf{Kata Kunci}: Rekrutmen, Onboarding, Kompensasi, Mediasi, BPJS}

\cventry{Agu 2012–Mei 2015}
        {Staf Rekrutmen}
        {PT Sumber Rejeki Abadi}
        {\url{https://www.sumberrejekiabadi.com}}
        {}
        {\begin{itemize}
\item Melakukan screening CV, wawancara awal, dan koordinasi dengan manajer departemen.
\item Mengelola database kandidat dan membuat laporan rekrutmen bulanan.
\item Mendukung program employer branding melalui job fair dan media sosial.
\end{itemize}
\textbf{Kata Kunci}: Screening, Wawancara, Employer Branding, Database}
\section{Bahasa}

\cvline{Indonesia}{Bahasa Ibu / Bilingual}
\cvline{Inggris}{Tingkat Profesional Penuh \hfill \textbf{Kata Kunci}: TOEFL 110, IELTS 8.0}
\cvline{Mandarin}{Tingkat Kerja Terbatas \hfill \textbf{Kata Kunci}: HSK 3}
\section{Keahlian}

\cvline{Manajemen Sumber Daya Manusia}{Ahli \hfill \textbf{Kata Kunci}: Rekrutmen, Manajemen Kinerja, Hubungan Karyawan, Kompensasi \& Benefit, Pelatihan \& Pengembangan}
\cvline{HRIS \& Analitik}{Mahir \hfill \textbf{Kata Kunci}: SAP SuccessFactors, Workday, Excel Lanjutan, Power BI, Tableau}
\cvline{Kepatuhan \& Hukum Ketenagakerjaan}{Mahir \hfill \textbf{Kata Kunci}: UU Ketenagakerjaan, BPJS, Perjanjian Kerja, Audit HR, Hubungan Industrial}
\cvline{Kepemimpinan \& Komunikasi}{Ahli \hfill \textbf{Kata Kunci}: Negosiasi, Mediasi, Public Speaking, Coaching, Manajemen Perubahan}
\section{Penghargaan}

\cventry{Nov 2022}
        {Asosiasi Manajemen Sumber Daya Manusia Indonesia}
        {HR Leader of the Year}
        {}
        {}
        {Diberikan kepada profesional HR yang menunjukkan kepemimpinan luar biasa dalam mengembangkan strategi SDM dan menciptakan budaya kerja positif.}

\cventry{Agu 2021}
        {HR Asia}
        {Employer Branding Terbaik}
        {}
        {}
        {Penghargaan atas keberhasilan program employer branding yang meningkatkan jumlah pelamar berkualitas dan reputasi perusahaan sebagai tempat kerja pilihan.}
\section{Sertifikat}

\cventry{Feb 2020}
        {Lembaga Sertifikasi Profesi HR Indonesia}
        {Certified Human Resources Professional (CHRP)}
        {\url{https://www.chrp-indonesia.org}}
        {}
        {}

\cventry{Jun 2018}
        {Society for Human Resource Management}
        {SHRM-CP}
        {\url{https://www.shrm.org/certification}}
        {}
        {}

\cventry{Mar 2017}
        {Kementerian Ketenagakerjaan RI}
        {Ahli K3 Umum}
        {\url{https://kemnaker.go.id}}
        {}
        {}
\section{Publikasi}

\cventry{Mei 2022}
        {Membangun Budaya Kerja Inklusif di Era Hybrid}
        {Jurnal Manajemen Indonesia}
        {\url{https://jurnalmanajemen.id/artikel/budaya-kerja-inklusif}}
        {}
        {\begin{itemize}
\item Artikel ini membahas strategi menciptakan budaya kerja inklusif bagi karyawan remote dan on-site.
\item Studi kasus pada tiga perusahaan teknologi di Jakarta menunjukkan peningkatan kolaborasi dan kepuasan kerja.
\end{itemize}}
\section{Referensi}

\cventry{\emaillink[siti.rahmawati@ui.ac.id]{siti.rahmawati@ui.ac.id}}
        {Dosen Psikologi Industri dan Organisasi}
        {Dr. Siti Rahmawati}
        {+62 811-9876-5432}
        {}
        {Budi Santoso adalah profesional HR yang berintegritas dan memiliki kemampuan analitis yang kuat. Ia mampu memimpin tim dan berkomunikasi efektif dengan berbagai pemangku kepentingan.}
\section{Proyek}

\cventry{Jan 2020–Des 2020}
        {Proyek transformasi digital untuk mengotomatisasi proses HR di perusahaan.}
        {Implementasi HRIS Terintegrasi}
        {}
        {}
        {\begin{itemize}
\item Memimpin migrasi data karyawan dari sistem manual ke platform HRIS berbasis cloud.
\item Melatih 50+ pengguna dan menyusun SOP baru untuk payroll, absensi, dan cuti.
\item Mencapai pengurangan waktu proses payroll dari 5 hari menjadi 2 hari.
\end{itemize}
\textbf{Kata Kunci}: HRIS, Digital Transformation, Change Management}

\cventry{Mar 2021–Agu 2021}
        {Inisiatif untuk meningkatkan kesejahteraan fisik dan mental karyawan.}
        {Program Kesejahteraan Karyawan}
        {}
        {}
        {\begin{itemize}
\item Meluncurkan program konseling online, kelas kebugaran, dan tunjangan kesehatan tambahan.
\item Mengukur dampak melalui survei kesejahteraan dan tingkat absensi.
\item Meningkatkan skor kesejahteraan karyawan sebesar 25\% dalam enam bulan.
\end{itemize}
\textbf{Kata Kunci}: Employee Wellbeing, Mental Health, Benefits}
\section{Minat}

\cvline{Pengembangan Diri}{Membaca, Public Speaking, Coaching}
\cvline{Olahraga}{Lari, Bersepeda, Bulu Tangkis}
\section{Kegiatan Sukarela}

\cventry{Jan 2019–Des 2022}
        {Mentor Karir}
        {Yayasan Pendidikan Anak Bangsa}
        {\url{https://www.ypab.org}}
        {}
        {\begin{itemize}
\item Membimbing siswa SMK dalam mempersiapkan diri memasuki dunia kerja.
\item Memberikan pelatihan CV, wawancara, dan keterampilan lunak.
\item Membantu menyelenggarakan job fair tahunan untuk 200+ siswa.
\end{itemize}}
    
\end{document}